% IN5450 Project 2 DOA Estimation
% Build: pdflatex -interaction=nonstopmode -halt-on-error presentation_p2.tex

\documentclass[aspectratio=169,9pt]{beamer}

\usetheme{Madrid}
\usecolortheme{default}

\usepackage{graphicx}
\usepackage{amsmath,amssymb}
\usepackage{booktabs}
\usepackage{mathtools}
\usepackage[T1]{fontenc}
\usepackage[utf8]{inputenc}
\usepackage{listings}

\graphicspath{{figs/}}

\setbeamertemplate{footline}[frame number]
\setbeamertemplate{navigation symbols}{}

% Tighter vertical spacing
\setlength{\parskip}{2pt plus 1pt}

% Figure helpers
\newcommand{\fullfig}[2][]{%
  \centering
  \includegraphics[width=\linewidth,height=0.75\textheight,keepaspectratio,#1]{#2}%
}
\newcommand{\colfig}[2][]{%
  \includegraphics[width=\linewidth,height=0.55\textheight,keepaspectratio,#1]{#2}%
}

% Listing style for MATLAB
\lstset{
  language=Matlab,
  basicstyle=\ttfamily\fontsize{5.5pt}{6.5pt}\selectfont,
  keywordstyle=\color{blue}\bfseries,
  commentstyle=\color{green!50!black},
  stringstyle=\color{red!70!black},
  breaklines=true,
  frame=single,
  numbers=none,
  columns=flexible
}

\title[IN5450 Project 2]{Direction of Arrival (DOA) Estimation Methods\\IN5450 Project 2}
\author{Theodor W\aa lberg}
\institute{University of Oslo}
\date{March 2026}

\begin{document}

% ----
\begin{frame}
  \titlepage
\end{frame}

% ----
\begin{frame}{Agenda}
  \tableofcontents
\end{frame}

% ==================================================================
\section{Signal Model \& Setup}
% ==================================================================

\begin{frame}{Signal model}
\vspace{-0.3em}
\textbf{Setup} (from \texttt{generate\_data.m}):
\begin{itemize}\setlength\itemsep{1pt}
  \item $M=10$ sensors, ULA with $d=0.5\lambda$
  \item $N=100$ snapshots
  \item Two incoherent sources: $\theta_1 = 0^\circ$, $\theta_2 = -10^\circ$
  \item Both at $\mathrm{SNR} = 0$\,dB, relative phase $\gamma = 45^\circ$
  \item Additive complex Gaussian noise ($\sigma^2 = 2$)
\end{itemize}

\medskip
\textbf{Data model:}
$$\mathbf{x}(t) = \mathbf{A}\,\mathbf{s}(t) + \mathbf{n}(t), \qquad
  \mathbf{A} = [\mathbf{a}(\theta_1)\;\;\mathbf{a}(\theta_2)], \qquad
  [\mathbf{a}(\theta)]_m = e^{-jmkd\sin\theta}$$

\medskip
\textbf{Sample correlation matrix:}
$$\hat{\mathbf{R}} = \frac{1}{N}\sum_{t=1}^{N}\mathbf{x}(t)\mathbf{x}^H(t) = \frac{1}{N}\mathbf{X}\mathbf{X}^H$$
\end{frame}

% ==================================================================
\section{Correlation Matrix}
% ==================================================================

\begin{frame}{Task 1 - Sample correlation matrix}
\vspace{-0.3em}
\fullfig{figure_01.pdf}

\vspace{-0.3em}
\scriptsize
$|\hat{\mathbf{R}}|$ shows the structure expected from a ULA.
\end{frame}

% ==================================================================
\section{Conventional Beamformer - DAS}
% ==================================================================

\begin{frame}{Task 2 - DAS spatial spectrum}
\vspace{-0.3em}
\begin{columns}[T]
\column{0.48\textwidth}
\colfig{figure_02.pdf}
\smallskip
\scriptsize \textbf{dB scale:} Single broad peak near $\theta\approx -5^\circ$.
The two sources at $0^\circ$ and $-10^\circ$ are \emph{not resolved}.

\column{0.48\textwidth}
\colfig{figure_03.pdf}
\smallskip
\scriptsize \textbf{Linear scale:} Confirms a single merged lobe.
\end{columns}

\medskip
\scriptsize
\textbf{Why unresolved?} The DAS beamwidth for $M\!=\!10$, $d\!=\!\lambda/2$ is $\approx\!11.5^\circ$,
which exceeds the $10^\circ$.
\end{frame}

% -- SLIDE: Code Task 2 --
\begin{frame}[fragile]{Code - DAS beamformer}
\vspace{-0.3em}
\begin{lstlisting}
R = x*x' / N;                     % Sample correlation matrix
d = 0.5;  kd = 2*pi*d;
theta = -40:0.25:50;

P_DAS = zeros(size(theta));
for n = 1:length(theta)
    DOA = theta(n);
    phi = -kd * sin(DOA * pi / 180);
    a   = exp(1i * phi) .^ (0:M-1)';
    P_DAS(n) = (a' * R * a) / M;
end
\end{lstlisting}

\smallskip
\scriptsize
$P_\text{DAS}(\theta) = \frac{1}{M}\mathbf{a}^H(\theta)\hat{\mathbf{R}}\,\mathbf{a}(\theta)$.
The conventional BF applies uniform weights. Resolution is limited by aperture size.
\end{frame}

% ==================================================================
\section{Capon / MVDR Beamformer}
% ==================================================================

\begin{frame}{Task 3 - Capon spatial spectrum}
\vspace{-0.3em}
\begin{columns}[T]
\column{0.48\textwidth}
\colfig{figure_04.pdf}
\smallskip
\scriptsize \textbf{dB scale:} Two distinct peaks visible near $0^\circ$ and $-10^\circ$.
The Capon BF resolves the sources that DAS could not.

\column{0.48\textwidth}
\colfig{figure_05.pdf}
\smallskip
\scriptsize \textbf{Linear scale:} Clear separation between the two source peaks.
\end{columns}

\medskip
\scriptsize
\textbf{Capon (MVDR):} $P_\text{Capon}(\theta) = \frac{1}{\mathbf{a}^H(\theta)\hat{\mathbf{R}}^{-1}\mathbf{a}(\theta)}$.
Minimises output power while maintaining unit gain in the look direction $\Rightarrow$ suppresses interference $\Rightarrow$ narrower peaks.
\end{frame}

% -- SLIDE: Code Task 3 --
\begin{frame}[fragile]{Code - Capon beamformer}
\vspace{-0.3em}
\begin{lstlisting}
Rinv = inv(R);
P_capon = zeros(size(theta));
for n = 1:length(theta)
    DOA = theta(n);
    phi = -kd * sin(DOA * pi / 180);
    a   = exp(1i * phi) .^ (0:M-1)';
    P_capon(n) = 1 / (a' * Rinv * a);
end
\end{lstlisting}

\smallskip
\scriptsize
Data-adaptive: uses $\hat{\mathbf{R}}^{-1}$ to place zeros on interferers.
Resolution exceeds DAS but is still finite.
\end{frame}

% ==================================================================
\section{Eigenvalue Decomposition}
% ==================================================================

\begin{frame}{Eigenvalue distribution of $\hat{\mathbf{R}}$}
\vspace{-0.3em}
\fullfig{figure_06.pdf}

\vspace{-0.3em}
\scriptsize
\textbf{Two dominant eigenvalues} $\gg$ remaining eight.

The gap between signal and noise eigenvalues enables subspace methods.
\end{frame}

% ==================================================================
\section{MUSIC Algorithm}
% ==================================================================

\begin{frame}{MUSIC spatial spectrum}
\vspace{-0.3em}
\begin{columns}[T]
\column{0.48\textwidth}
\colfig{figure_07.pdf}
\smallskip
\scriptsize \textbf{dB scale:} Sharp peaks at $\theta_1\!=\!0^\circ$ and $\theta_2\!=\!-10^\circ$.
MUSIC achieves super-resolution.

\column{0.48\textwidth}
\colfig{figure_08.pdf}
\smallskip
\scriptsize \textbf{Linear scale:} Both sources clearly resolved with very narrow peaks.
\end{columns}

\medskip
\scriptsize
$P_\text{MUSIC}(\theta) = \frac{\mathbf{a}^H\mathbf{a}}{\mathbf{a}^H\mathbf{U}_n\mathbf{U}_n^H\mathbf{a}}$.
MUSIC exploits orthogonality between signal steering vectors and noise subspace:
$\mathbf{U}_n^H\mathbf{a}(\theta_k) = \mathbf{0}$ at true DOAs $\Rightarrow$ denominator $\to 0$ $\Rightarrow$ sharp peaks.
\end{frame}

% -- SLIDE: Code Task 5 --
\begin{frame}[fragile]{Code - MUSIC}
\vspace{-0.3em}
\begin{lstlisting}
[V,D] = eig(R);
[dd,I] = sort(diag(D));
dd = flipud(dd);  V = V(:,flipud(I));  % descending order

num_signals = 2;
Un = V(:, num_signals+1 : end);   % noise subspace
Pn = Un * Un';

P_music = zeros(size(theta));
for n = 1:length(theta)
    DOA = theta(n);
    phi = -kd * sin(DOA * pi / 180);
    a   = exp(1i * phi) .^ (0:M-1)';
    P_music(n) = (a' * a) / (a' * Pn * a);
end
\end{lstlisting}

\smallskip
\scriptsize
Signal subspace: eigenvectors for the $K\!=\!2$ largest eigenvalues.
Noise subspace: remaining $M\!-\!K\!=\!8$ eigenvectors. The spectrum peaks where $\mathbf{a}(\theta)$ is orthogonal to $\mathbf{U}_n$.
\end{frame}

% ==================================================================
\section{Eigenvector Method}
% ==================================================================

\begin{frame}{Task 6 - Eigenvector method}
\vspace{-0.3em}
\begin{columns}[T]
\column{0.48\textwidth}
\colfig{figure_09.pdf}
\smallskip
\scriptsize \textbf{dB scale:} Peaks at $0^\circ$ and $-10^\circ$, similar to MUSIC.

\column{0.48\textwidth}
\colfig{figure_10.pdf}
\smallskip
\scriptsize \textbf{Linear scale:} Peak structure comparable to MUSIC with slight amplitude differences.
\end{columns}

\medskip
\scriptsize
\textbf{Eigenvector method:}
$P_\text{EV}(\theta) = \frac{1}{\mathbf{a}^H(\theta)\hat{\mathbf{R}}_n^{-1}\mathbf{a}(\theta)}$,
where $\hat{\mathbf{R}}_n^{-1} = \sum_{i=K+1}^{M}\frac{1}{\lambda_i}\mathbf{v}_i\mathbf{v}_i^H$.

Compared to MUSIC, the EV method weights each noise eigenvector by $1/\lambda_i$,
giving more weight to eigenvectors with smaller eigenvalues, which can sharpen peaks.
\end{frame}

% -- SLIDE: Code Task 6 --
\begin{frame}[fragile]{Code - Eigenvector method}
\vspace{-0.3em}
\begin{lstlisting}
Vn = V(:, num_signals+1:end);
Ln = dd(num_signals+1:end);
R_noise_only_inv = Vn * diag(1 ./ Ln) * Vn';

P_EV = zeros(size(theta));
for n = 1:length(theta)
    DOA = theta(n);
    phi = -kd * sin(DOA * pi / 180);
    a   = exp(1i*phi).^(0:M-1)';
    P_EV(n) = 1 / (a' * R_noise_only_inv * a);
end
\end{lstlisting}

\smallskip
\scriptsize
Key difference from MUSIC: each noise eigenvector weighted by $1/\lambda_i$ instead of unity.
Both methods rely on correct estimation of $K$ number of sources.
\end{frame}

% ==================================================================
\section{Incorrect Source Count}
% ==================================================================

\begin{frame}{Task 7 - Effect of incorrect $\hat{K}$ on MUSIC \& EV}
\vspace{-0.3em}
$\hat{K}=0$ (no signal subspace. Noise subspace is the full space):
\fullfig[height=0.60\textheight]{figure_11.pdf}

\vspace{-0.3em}
\scriptsize
With $\hat{K}\!=\!0$ the ``noise subspace'' spans all $M$ dimensions, so $\mathbf{U}_n\mathbf{U}_n^H = \mathbf{I}$.
MUSIC spectrum becomes $P = M/\|\mathbf{a}\|^2 = 1$ flat, no peaks. EV spectrum is $1/(\mathbf{a}^H\mathbf{R}^{-1}\mathbf{a})$ = Capon.
\end{frame}

\begin{frame}{Task 7 - $\hat{K}=1$}
\vspace{-0.3em}
\fullfig[height=0.65\textheight]{figure_12.pdf}

\vspace{-0.3em}
\scriptsize
With $\hat{K}\!=\!1$: one signal eigenvector is misclassified as noise.
Only one peak is visible; the algorithm cannot resolve both sources.
The noise subspace is ``contaminated'' by signal energy $\Rightarrow$ the orthogonality condition
$\mathbf{U}_n^H\mathbf{a}(\theta_k)=\mathbf{0}$ no longer holds for both DOAs.
\end{frame}

\begin{frame}{Task 7 - $\hat{K}=3$}
\vspace{-0.3em}
\fullfig[height=0.65\textheight]{figure_13.pdf}

\vspace{-0.3em}
\scriptsize
With $\hat{K}\!=\!3$: the noise subspace loses one dimension.
Both true peaks remain, but a spurious third peak may appear,
and the noise floor rises because the noise subspace is under represented.
\end{frame}

% -- SLIDE: Code Task 7 --
\begin{frame}[fragile]{Code - Incorrect source count}
\vspace{-0.3em}
\begin{lstlisting}
Ns_estimates = [0, 1, 3];

for k = 1:length(Ns_estimates)
    Ns_est = Ns_estimates(k);
    Vn = V(:, Ns_est+1 : end);
    Ln = dd(Ns_est+1 : end);

    for n = 1:length(theta)
        DOA = theta(n);
        phi = -kd * sin(DOA * pi / 180);
        a   = exp(1i * phi) .^ (0:M-1)';

        % MUSIC
        P_music_task7(n) = 1 / (a' * (Vn * Vn') * a);

        % Eigenvector
        P_EV_task7(n) = 1 / (a' * (Vn * diag(1./Ln) * Vn') * a);
    end
end
\end{lstlisting}

\smallskip
\scriptsize
Correct $\hat{K}$ is critical: underestimation loses sources, overestimation raises the noise floor
and can introduce false peaks.
\end{frame}

% ==================================================================
\section*{Summary}
% ==================================================================

\begin{frame}{Summary}
\begin{enumerate}\setlength\itemsep{3pt}
  \item \textbf{Correlation matrix (1):} Structure expected from a ULA; diagonal $\approx \sigma_s^2 + \sigma_n^2$.
  \item \textbf{DAS (2):} Cannot resolve sources.
  \item \textbf{Capon (3):} Data-adaptive; resolves the two sources by suppressing interference.
  \item \textbf{Eigenvalues (4):} Two dominant eigenvalues confirm $K\!=\!2$ sources; gap separates signal and noise subspaces.
  \item \textbf{MUSIC (5):} Super-resolution via noise-subspace orthogonality; sharp peaks at true DOAs.
  \item \textbf{Eigenvector (6):} Eigenvalue-weighted noise subspace; similar resolution to MUSIC.
  \item \textbf{Incorrect $\hat{K}$ (7):} Underestimation merges/loses peaks; overestimation raises noise floor.
\end{enumerate}
\end{frame}

\end{document}
